\section{Related Work}
\label{sec:related}

\textbf{Model Merging and Weight Averaging.} Model merging has emerged as an active area of research to consolidate specialized models into a single multi-task network without explicit multi-task training. Early iterations such as Model Soups \cite{modelsoups} and SWA \cite{weight_averaging} demonstrated that simple linear averaging of weights can improve generalization when models are fine-tuned from a shared pre-trained initialization. Later, Task Arithmetic \cite{taskarithmetic} extended this concept by showing that fine-tuned updates (task vectors) can be linearly combined (added or subtracted) to modulate task-specific behaviors. However, interference between parameter updates can lead to performance degradation. To address parameter conflict, TIES-Merging \cite{tiesmerging} introduces a sign-consensus protocol and prunes redundant updates, while DARE \cite{dare} utilizes random sparsification and rescaling. Other works explore permutation alignment to resolve coordinate symmetries before merging, such as Git Re-basin \cite{gitrebasin}, ZipIt \cite{zipit}, and REASON \cite{reason}. In contrast to these post-hoc weight transformations, we evaluate whether complex post-merge manipulations are rendered redundant by proper optimization.

\textbf{Manifold and Spectral Merging.} Rather than performing Euclidean averaging, several recent works propose mapping or constraining model updates to non-Euclidean manifolds or spectral subspaces. For instance, OrthoMerge \cite{orthomerge2026} maps orthogonal weight layers to Lie algebra to merge models on the Riemannian manifold of orthogonal groups. Similarly, SAIM \cite{saim2026} proposes that individual model updates undergo an SVD-based adaptive isotropic merging step, which interpolates the singular values of combined weight matrices toward their mean to balance the spectrum and prevent representation collapse. Other works suggest merging models in concrete subspaces \cite{concrete_subspace} or using Fisher-weighted averaging \cite{fisher_merging}. Our methodological critique specifically investigates SAIM's spectral interpolation, proving that such SVD-based reconstructions can actually disrupt task-specific representations and degrade performance.

\textbf{Sharpness-Aware Minimization and Flatness.} Modern deep learning theory suggests that the geometry of the loss landscape around a local minimum is strongly correlated with a model's generalization capabilities \cite{weight_averaging, sam}. Sharpness-Aware Minimization (SAM) \cite{sam} formalizes this by explicitly minimizing the worst-case neighborhood loss, finding wider and flatter minima. Extensions of SAM include scale-invariant formulations (ASAM) \cite{asampaper}, gradient-scaled perturbations (GSAM) \cite{gsampaper}, and coordinate-restricted sharpness minimization (CSAM) \cite{csampaper}. In the context of model merging, LMC (Linear Mode Connectivity) studies have shown that models residing in flat basins can be interpolated without encountering high-loss barriers \cite{lmc1, lmc2}. We build on this insight to demonstrate that optimizer-driven flatness is the principal causal driver of model merging success, making complex post-hoc spectral transformations unnecessary.

\textbf{Continual Learning via Sequential Merging.} Traditional continual learning relies on parameter regularization \cite{ewc, mas, si}, rehearsal buffers \cite{gem, der}, or dynamic architectures \cite{gpm, hat}. More recently, sequential model merging has been proposed as a training-free and buffer-free alternative \cite{clmerging}. SyMerge \cite{symerge2026} proposes adapting a single task-specific layer of a transformer and optimizing merging coefficients to achieve task synergy. SAIM \cite{saim2026} combines its SA-BCD optimizer with isotropic SVD merging to enable robust sequential merging over a long horizon of tasks. In this paper, we audit these claims, demonstrating that sequential merging is indeed viable, but can be executed far more effectively and with less computational overhead by simply combining standard SAM with naive Task Arithmetic.
